\documentclass{article}
\usepackage[T1]{fontenc}
\usepackage{lmodern}
\usepackage{graphicx}
\usepackage{amsmath,amssymb}
\usepackage[a4paper,margin=2cm]{geometry}
\usepackage[absolute,overlay]{textpos}
\setlength{\TPHorizModule}{1mm}
\setlength{\TPVertModule}{1mm}
\usepackage[indonesian]{babel}
\usepackage{hyperref}
\setlength{\emergencystretch}{2em}
% unit: Halaman 15

\begin{document}

% === Halaman 15 (llm) ===

\textbf{Penggunaan kain sutra dan lilin}

\vspace{121.77mm}

\begin{itemize}
    \item Orang Cina kuno menulis catatan pada potongan-potongan kecil sutra yang kemudian digumpalkan menjadi bola kecil dan dilapisi lilin.
    \vspace{20mm}
    \item Selanjutnya bola kecil tersebut ditelan oleh si pembawa pesan.
    \item Pesan dibaca setelah bola kecil dikeluarkan dari perut si pembawa pesan dengan cara BAB.
\end{itemize}

% deskripsi: Ilustrasi kartun seekor panda yang sedang memegang daun hijau.
% bbox normalisasi: [0.7100, 0.3300, 0.8900, 0.7400]
\begin{textblock*}{37.80mm}(149.10mm,98.01mm)
\noindent\includegraphics[width=37.80mm,height=121.77mm]{../../images/llm/page_015_img_01.png}%
\end{textblock*}

\end{document}
